\documentclass[10pt,a4paper]{amsart}
\usepackage{amsmath,amssymb,amsthm,graphicx}
% Retained minimal preamble for the Stacks Project chapter sites-cohomology.tex (Cohomology on Sites).
% Verbatim excerpt of the chapter preamble: Xy-pic setup (native preamble lines 14-18) and the native
% theorem declarations with the mathematical macros (native preamble lines 154-217). Copied mechanically;
% nothing added, removed or redefined. The wrapper class and the amsmath/amssymb/amsthm/graphicx base
% packages are supplied by the assembler, so the class/externaldocument part of the native preamble is omitted.
% For commutative diagrams we use Xy-pic
\usepackage[all]{xy}

% We use 2cell for 2-commutative diagrams.
\xyoption{2cell}
% Theorem environments.
%
\theoremstyle{plain}
\newtheorem{theorem}[subsection]{Theorem}
\newtheorem{proposition}[subsection]{Proposition}
\newtheorem{lemma}[subsection]{Lemma}

\theoremstyle{definition}
\newtheorem{definition}[subsection]{Definition}
\newtheorem{example}[subsection]{Example}
\newtheorem{exercise}[subsection]{Exercise}
\newtheorem{situation}[subsection]{Situation}

\theoremstyle{remark}
\newtheorem{remark}[subsection]{Remark}
\newtheorem{remarks}[subsection]{Remarks}

\numberwithin{equation}{subsection}

% Macros
%
\def\lim{\mathop{\mathrm{lim}}\nolimits}
\def\colim{\mathop{\mathrm{colim}}\nolimits}
\def\Spec{\mathop{\mathrm{Spec}}}
\def\Hom{\mathop{\mathrm{Hom}}\nolimits}
\def\Ext{\mathop{\mathrm{Ext}}\nolimits}
\def\SheafHom{\mathop{\mathcal{H}\!\mathit{om}}\nolimits}
\def\SheafExt{\mathop{\mathcal{E}\!\mathit{xt}}\nolimits}
\def\Sch{\mathit{Sch}}
\def\Mor{\mathop{\mathrm{Mor}}\nolimits}
\def\Ob{\mathop{\mathrm{Ob}}\nolimits}
\def\Sh{\mathop{\mathit{Sh}}\nolimits}
\def\NL{\mathop{N\!L}\nolimits}
\def\CH{\mathop{\mathrm{CH}}\nolimits}
\def\proetale{{pro\text{-}\acute{e}tale}}
\def\etale{{\acute{e}tale}}
\def\QCoh{\mathit{QCoh}}
\def\Ker{\mathop{\mathrm{Ker}}}
\def\Im{\mathop{\mathrm{Im}}}
\def\Coker{\mathop{\mathrm{Coker}}}
\def\Coim{\mathop{\mathrm{Coim}}}

% Boxtimes
%
\DeclareMathSymbol{\boxtimes}{\mathbin}{AMSa}{"02}

%
% Macros for moduli stacks/spaces
%
\def\QCohstack{\mathcal{QC}\!\mathit{oh}}
\def\Cohstack{\mathcal{C}\!\mathit{oh}}
\def\Spacesstack{\mathcal{S}\!\mathit{paces}}
\def\Quotfunctor{\mathrm{Quot}}
\def\Hilbfunctor{\mathrm{Hilb}}
\def\Curvesstack{\mathcal{C}\!\mathit{urves}}
\def\Polarizedstack{\mathcal{P}\!\mathit{olarized}}
\def\Complexesstack{\mathcal{C}\!\mathit{omplexes}}
% \Pic is the operator that assigns to X its picard group, usage \Pic(X)
% \Picardstack_{X/B} denotes the Picard stack of X over B
% \Picardfunctor_{X/B} denotes the Picard functor of X over B
\def\Pic{\mathop{\mathrm{Pic}}\nolimits}
\def\Picardstack{\mathcal{P}\!\mathit{ic}}
\def\Picardfunctor{\mathrm{Pic}}
\def\Deformationcategory{\mathcal{D}\!\mathit{ef}}

% --- appended 2026-10-12 (Chief): the project own definitions of the reference/slogan/history
% display environments, copied mechanically from the shared Stacks preamble lines 8-12
% (\usepackage{verbatim} plus the three \newenvironment lines, which the chapter preamble keeps
% verbatim from the project preamble). Nothing added, removed or redefined: these environments
% carry no mathematical content and the project itself typesets them as comments. Without these
% three lines a counted statement whose own text uses them fails with Environment slogan undefined.
\usepackage{verbatim}
\newenvironment{reference}{\comment}{\endcomment}
\newenvironment{slogan}{\comment}{\endcomment}
\newenvironment{history}{\comment}{\endcomment}
\usepackage{cleveref}
\crefname{originallemma}{original lemma}{original lemmas}
\Crefname{originallemma}{Original lemma}{Original lemmas}
\makeatletter
\newlabel{homology-lemma-five-lemma}{{05QB}{source}{}{sourcepointer-1}{}}
\newlabel{homology-lemma-five-lemma@cref}{{[originallemma][0][]05QB}{[1][0][]source}{}{}{}}
\makeatother
\crefname{originallemma}{original lemma}{original lemmas}
\Crefname{originallemma}{Original lemma}{Original lemmas}
\makeatletter
\newlabel{homology-lemma-four-lemma}{{05QA}{source}{}{sourcepointer-2}{}}
\newlabel{homology-lemma-four-lemma@cref}{{[originallemma][0][]05QA}{[1][0][]source}{}{}{}}
\makeatother
\crefname{originalreference}{original reference}{original references}
\Crefname{originalreference}{Original reference}{Original references}
\makeatletter
\newlabel{lemma-cone}{{08C5}{source}{}{sourcepointer-3}{}}
\newlabel{lemma-cone@cref}{{[originalreference][0][]08C5}{[1][0][]source}{}{}{}}
\makeatother
\crefname{originalreference}{original reference}{original references}
\Crefname{originalreference}{Original reference}{Original references}
\makeatletter
\newlabel{lemma-lift-through-quasi-isomorphism}{{08C8}{source}{}{sourcepointer-4}{}}
\newlabel{lemma-lift-through-quasi-isomorphism@cref}{{[originalreference][0][]08C8}{[1][0][]source}{}{}{}}
\makeatother
\crefname{originalreference}{original reference}{original references}
\Crefname{originalreference}{Original reference}{Original references}
\makeatletter
\newlabel{lemma-pseudo-coherent-independent-representative}{{08CC}{source}{}{sourcepointer-5}{}}
\newlabel{lemma-pseudo-coherent-independent-representative@cref}{{[originalreference][0][]08CC}{[1][0][]source}{}{}{}}
\makeatother
\title{Pseudo-coherence in a distinguished triangle, with one degree of slack}
\author{The Stacks Project Authors}
\date{stacks-project:sites-cohomology.tex:local-snapshot-20261011; GNU FDL 1.2 or later}
\begin{document}
\maketitle
\noindent\textit{Editorial note.} All delivered mathematical text is the author's own, copied from the pinned Stacks Project chapter file sites-cohomology.tex; no statement or proof step was rewritten, reordered or removed.\par\medskip
\noindent\textit{Editorial note.} The counted proof is the author's complete argument: the proof is written out in full, with the author's own remark that some details of representing the maps by complexes are omitted; the construction of the cone, the morphism of triangles and the comparison of the long exact sequences are all given.\par\medskip
\noindent\textit{Editorial note.} The label lemma-cone-pseudo-coherent names the counted statement, and the proof environment immediately following it in the source is delivered verbatim as its proof.\par\medskip
\noindent\textit{Editorial note.} The counted statement is carried by the author's own lemma environment in the chapter source.\par\medskip
\noindent\textit{Editorial note.} This card belongs to the priority queue of the D-72 twin ruling (the sites-side statement is the designated counting side). The statement and proof are copied from the pinned chapter file; the five cited inputs are kept as external pointers with their canonical Stacks tags.\par\medskip
\noindent\textit{Allowed standard prerequisites (editorial).}\begin{itemize}
\item Lifting a map of complexes through a quasi-isomorphism. Same chapter, the lemma allowing a map in the derived category to be represented by a map of complexes after refining a covering. Supplies the maps gamma\_i of complexes used to build the cone.
\item Strict perfectness of the cone on a morphism of strictly perfect complexes. Same chapter, the lemma stating that the cone on a morphism of strictly perfect complexes is strictly perfect. Supplies strict perfectness of C(gamma\_i).
\item The four and five lemmas on long exact cohomology sequences, and locality of pseudo-coherence. Homology, the four and five lemmas, and Same chapter, the locality lemma for pseudo-coherence. Supplies the conclusion that the induced map on the cone has the required cohomological behaviour, hence that M is m-pseudo-coherent.
\end{itemize}\medskip
\ifdefined\hypertarget\hypertarget{sourcepointer-1}{}\fi
\noindent\textit{Source reference.} Original lemma 05QB: sources/stacks-chow/tags.txt, source lines 7370--7370. Prerequisite pointer invoked at proof line(s) 12226; it names a standard fact of a sibling chapter that the counted argument uses, and every step of the counted proof built on it is retained verbatim in the delivered body.\par\medskip
\ifdefined\hypertarget\hypertarget{sourcepointer-2}{}\fi
\noindent\textit{Source reference.} Original lemma 05QA: sources/stacks-chow/tags.txt, source lines 7369--7369. Prerequisite pointer invoked at proof line(s) 12225; it names a standard fact of a sibling chapter that the counted argument uses, and every step of the counted proof built on it is retained verbatim in the delivered body.\par\medskip
\ifdefined\hypertarget\hypertarget{sourcepointer-3}{}\fi
\noindent\textit{Source reference.} Original reference 08C5: sources/stacks-sites-cohomology/pinned-original-source.tex, source lines 11735--11735. Prerequisite pointer invoked at proof line(s) 12216; it names a standard fact of the same chapter that the counted argument uses, and every step of the counted proof built on it is retained verbatim in the delivered body.\par\medskip
\ifdefined\hypertarget\hypertarget{sourcepointer-4}{}\fi
\noindent\textit{Source reference.} Original reference 08C8: sources/stacks-sites-cohomology/pinned-original-source.tex, source lines 11848--11848. Prerequisite pointer invoked at proof line(s) 12202; it names a standard fact of the same chapter that the counted argument uses, and every step of the counted proof built on it is retained verbatim in the delivered body.\par\medskip
\ifdefined\hypertarget\hypertarget{sourcepointer-5}{}\fi
\noindent\textit{Source reference.} Original reference 08CC: sources/stacks-sites-cohomology/pinned-original-source.tex, source lines 12066--12066. Prerequisite pointer invoked at proof line(s) 12230; it names a standard fact of the same chapter that the counted argument uses, and every step of the counted proof built on it is retained verbatim in the delivered body.\par\medskip
\begin{lemma}
\label{lemma-cone-pseudo-coherent}
Let $(\mathcal{C}, \mathcal{O})$ be a ringed site and $m \in \mathbf{Z}$.
Let $(K, L, M, f, g, h)$ be a distinguished triangle in $D(\mathcal{O})$.
\begin{enumerate}
\item If $K$ is $(m + 1)$-pseudo-coherent and $L$ is $m$-pseudo-coherent
then $M$ is $m$-pseudo-coherent.
\item If $K$ and $M$ are $m$-pseudo-coherent, then $L$ is $m$-pseudo-coherent.
\item If $L$ is $(m + 1)$-pseudo-coherent and $M$
is $m$-pseudo-coherent, then $K$ is $(m + 1)$-pseudo-coherent.
\end{enumerate}
\end{lemma}

\begin{proof}
Proof of (1). Let $U$ be an object of $\mathcal{C}$. Choose a covering
$\{U_i \to U\}$ and maps $\alpha_i : \mathcal{K}_i^\bullet \to K|_{U_i}$
in $D(\mathcal{O}_{U_i})$ with $\mathcal{K}_i^\bullet$ strictly perfect and
$H^j(\alpha_i)$ isomorphisms for $j > m + 1$ and surjective for $j = m + 1$.
We may replace $\mathcal{K}_i^\bullet$ by
$\sigma_{\geq m + 1}\mathcal{K}_i^\bullet$
and hence we may assume that $\mathcal{K}_i^j = 0$
for $j < m + 1$. After refining the covering we may choose
maps $\beta_i : \mathcal{L}_i^\bullet \to L|_{U_i}$ in $D(\mathcal{O}_{U_i})$
with $\mathcal{L}_i^\bullet$ strictly perfect such that
$H^j(\beta)$ is an isomorphism for $j > m$ and
surjective for $j = m$. By
Lemma \ref{lemma-lift-through-quasi-isomorphism}
we can, after refining the covering, find maps of complexes
$\gamma_i : \mathcal{K}^\bullet \to \mathcal{L}^\bullet$
such that the diagrams
$$
\xymatrix{
K|_{U_i} \ar[r] & L|_{U_i} \\
\mathcal{K}_i^\bullet \ar[u]^{\alpha_i} \ar[r]^{\gamma_i} &
\mathcal{L}_i^\bullet \ar[u]_{\beta_i}
}
$$
are commutative in $D(\mathcal{O}_{U_i})$ (this requires representing the
maps $\alpha_i$, $\beta_i$ and $K|_{U_i} \to L|_{U_i}$
by actual maps of complexes; some details omitted).
The cone $C(\gamma_i)^\bullet$ is strictly perfect (Lemma \ref{lemma-cone}).
The commutativity of the diagram implies that there exists a morphism
of distinguished triangles
$$
(\mathcal{K}_i^\bullet, \mathcal{L}_i^\bullet, C(\gamma_i)^\bullet)
\longrightarrow
(K|_{U_i}, L|_{U_i}, M|_{U_i}).
$$
It follows from the induced map on long exact cohomology sequences and
Homology, Lemmas \ref{homology-lemma-four-lemma} and
\ref{homology-lemma-five-lemma}
that $C(\gamma_i)^\bullet \to M|_{U_i}$ induces an isomorphism
on cohomology in degrees $> m$ and a surjection in degree $m$.
Hence $M$ is $m$-pseudo-coherent by
Lemma \ref{lemma-pseudo-coherent-independent-representative}.

\medskip\noindent
Assertions (2) and (3) follow from (1) by rotating the distinguished
triangle.
\end{proof}

\end{document}
